\section{Conclusion}\label{sec:conclusion}

We built Task Manager, a JavaFX desktop application over Spring Boot services that adds, lists, updates and deletes
tasks, sets their status, refuses a task without a title, and keeps its tasks in a local SQLite file. The rules live
in one place, the services, where the JUnit tests check them on a fresh database, and the window only shows what the
services return or refuse. All 38 tests pass, the Javadoc HTML is generated with every check on and not one warning,
and the history records the work as \CommitCount\ commits, each with a message that says what changed and why.

The lab's other half was the workflow. Working in an IDE on the JDK meant running and debugging the program where it
was written. Committing each piece of work on its own, after reading its diff and running the tests, left a history
that explains how the program was built. Writing a Javadoc comment for every member, private ones included, made us
state what each one is for, and the generated pages show the design as a reader meets it.

Not everything worked the first time. A title over its limit at first reached the database's own checks instead of a
message, and a save could be overtaken by a delete in a second copy of the program; both were found, reproduced by a
test and fixed, as Section~\ref{sec:git} shows. Two limits remain, both deliberate for a single-user program: when two
copies edit the same task that still exists, the last save wins without a warning, and the search ignores case only
for the letters A to Z. An optimistic lock on the task, and a search that folds the case of every letter, would be the
next steps.
